\documentclass{beamer}
\usepackage[utf8]{inputenc}
\usepackage[T1]{fontenc}
\usepackage[spanish]{babel}
\usepackage{graphicx}
\usepackage{tikz}
\usetikzlibrary{shapes,arrows,positioning}

% Tema
\usetheme{Madrid}
\usecolortheme{default}

% Colores personalizados UNS
\definecolor{unsazul}{RGB}{30,60,120}
\setbeamercolor{structure}{fg=unsazul}
\setbeamercolor{palette primary}{bg=unsazul,fg=white}
\setbeamercolor{palette secondary}{bg=unsazul!80,fg=white}
\setbeamercolor{palette tertiary}{bg=unsazul!60,fg=white}

% Información del documento
\title{Métodos de Análisis de Datos 1}
\subtitle{Clase 1: El Ciclo de Análisis de Datos}
\author{Universidad Nacional del Sur}
\institute{}
\date{}

\begin{document}

% Portada
\begin{frame}
\titlepage
\end{frame}

% Contenidos
\begin{frame}{Contenidos de hoy}
\tableofcontents
\end{frame}

\section{Información del curso}

\begin{frame}{Métodos de Análisis de Datos 1}
\begin{block}{Información del curso}
\begin{itemize}
    \item \textbf{Modalidad:} Presencial
    \item \textbf{Horario:} Lunes y miércoles
    \item \textbf{Duración:} 30 clases (16 marzo - 3 julio)
    \item \textbf{Herramientas:} Python, R, VS Code, Jupyter
\end{itemize}
\end{block}

\begin{block}{Evaluación}
\begin{itemize}
    \item Sin parciales escritos
    \item Entrega de Proyecto integrador al final del curso
    \item Github personal para subir tareas y proyecto
\end{itemize}
\end{block}
\end{frame}

\begin{frame}{Objetivos del curso}
Al finalizar el curso serán capaces de:

\begin{itemize}
    \item Aplicar el \textbf{ciclo completo} de análisis de datos
    \item Utilizar \textbf{Python y R} para análisis de datos
    \item Construir y evaluar \textbf{modelos} estadísticos y de ML
    \item \textbf{Comunicar} hallazgos mediante visualizaciones y reportes
    \item Trabajar con \textbf{datos reales} de diversas fuentes
\end{itemize}
\end{frame}

\section{¿Qué es el análisis de datos?}

\begin{frame}{¿Qué es el análisis de datos?}
\begin{block}{Definición}
El \textbf{análisis de datos} es el proceso de inspeccionar, limpiar, transformar y modelar datos con el objetivo de:
\begin{itemize}
    \item Descubrir información útil
    \item Informar conclusiones
    \item Apoyar la toma de decisiones
\end{itemize}
\end{block}
\end{frame}

\begin{frame}{¿Qué es el análisis de datos?}
\begin{block}{Definición}
El \textbf{análisis de datos} es el proceso de inspeccionar, limpiar, transformar y modelar datos con el objetivo de:
\begin{itemize}
    \item Descubrir información útil
    \item Informar conclusiones
    \item Apoyar la toma de decisiones
\end{itemize}
\end{block}

\begin{exampleblock}{Ejemplo}
Una empresa de e-commerce quiere entender por qué bajaron las ventas:
\begin{itemize}
    \item ¿Qué productos tienen mayor abandono del carrito?
    \item ¿Qué clientes tienen mayor probabilidad de comprar?
    \item ¿Qué acciones aumentarían las ventas?
\end{itemize}
\end{exampleblock}
\end{frame}

\section{El ciclo de análisis de datos}

\begin{frame}{El ciclo de análisis de datos}
\begin{center}
\begin{tikzpicture}[scale=0.8, every node/.style={scale=0.8}]
    % Círculo central
    \draw[thick] (0,0) circle (3cm);
    
    % Etapas (6 puntos en el círculo)
    \foreach \angle/\num/\label in {
        90/1/Obtencion,
        30/2/Limpieza,
        -30/3/Exploracion,
        -90/4/Modelado,
        -150/5/Prediccion,
        150/6/Comunicacion
    } {
        \node[circle, fill=unsazul, text=white, minimum size=0.8cm] at (\angle:3cm) {\small\num};
        \node[align=center] at (\angle:3.8cm) {\small\textbf{\label}};
    }
    
    % Flechas circulares
    \foreach \start/\toAngle in {90/30, 30/-30, -30/-90, -90/-150, -150/-210,-210/-269 } {
        \draw[->, thick, unsazul] (\start:2.5cm) arc (\start:\toAngle:2.5cm);
    }
\end{tikzpicture}
\end{center}

\vspace{0.3cm}

El ciclo es iterativo: podemos volver a etapas anteriores
\end{frame}



\begin{frame}{Etapas del ciclo}
\begin{enumerate}
    \item \textbf{Obtención:} Acceder a datos (archivos, APIs, bases de datos, web scraping)
\end{enumerate}
\end{frame}

\begin{frame}{Etapas del ciclo}
\begin{enumerate}
    \item \textbf{Obtención:} Acceder a datos (archivos, APIs, bases de datos, web scraping)
    \item \textbf{Limpieza:} Corregir errores, tratar valores faltantes, transformar
\end{enumerate}
\end{frame}

\begin{frame}{Etapas del ciclo}
\begin{enumerate}
    \item \textbf{Obtención:} Acceder a datos (archivos, APIs, bases de datos, web scraping)
    \item \textbf{Limpieza:} Corregir errores, tratar valores faltantes, transformar
    \item \textbf{Exploración:} Estadística descriptiva y visualizaciones
\end{enumerate}
\end{frame}

\begin{frame}{Etapas del ciclo}
\begin{enumerate}
    \item \textbf{Obtención:} Acceder a datos (archivos, APIs, bases de datos, web scraping)
    \item \textbf{Limpieza:} Corregir errores, tratar valores faltantes, transformar
    \item \textbf{Exploración:} Estadística descriptiva y visualizaciones
    \item \textbf{Modelado:} Construir modelos para responder preguntas
\end{enumerate}
\end{frame}

\begin{frame}{Etapas del ciclo}
\begin{enumerate}
    \item \textbf{Obtención:} Acceder a datos (archivos, APIs, bases de datos, web scraping)
    \item \textbf{Limpieza:} Corregir errores, tratar valores faltantes, transformar
    \item \textbf{Exploración:} Estadística descriptiva y visualizaciones
    \item \textbf{Modelado:} Construir modelos para responder preguntas
    \item \textbf{Predicción:} Aplicar modelos a nuevos datos
\end{enumerate}
\end{frame}

\begin{frame}{Etapas del ciclo}
\begin{enumerate}
    \item \textbf{Obtención:} Acceder a datos (archivos, APIs, bases de datos, web scraping)
    \item \textbf{Limpieza:} Corregir errores, tratar valores faltantes, transformar
    \item \textbf{Exploración:} Estadística descriptiva y visualizaciones
    \item \textbf{Modelado:} Construir modelos para responder preguntas
    \item \textbf{Predicción:} Aplicar modelos a nuevos datos
    \item \textbf{Comunicación:} Presentar hallazgos de forma clara
\end{enumerate}
\end{frame}

\section{Tipos de análisis}

\begin{frame}{Tipos de análisis de datos}
\begin{columns}[t]
\column{0.48\textwidth}
\begin{block}{Descriptivo}
\textbf{¿Qué pasó?}\\
Estadísticas resumidas y visualizaciones\\
\textit{Ej: ``Las ventas fueron \$50,000''}
\end{block}

\begin{block}{Exploratorio (EDA)}
\textbf{¿Qué patrones hay?}\\
Buscar relaciones y tendencias\\
\textit{Ej: ``Ventas bajan los fines de semana''}
\end{block}

\column{0.48\textwidth}
\begin{block}{Inferencial}
\textbf{¿Podemos generalizar?}\\
Inferencias con incertidumbre\\
\textit{Ej: ``El descuento aumenta ventas 10-25\%''}
\end{block}

\begin{block}{Predictivo}
\textbf{¿Qué pasará?}\\
Modelos para el futuro\\
\textit{Ej: ``Cliente X comprará con 78\% prob.''}
\end{block}
\end{columns}
\end{frame}

\begin{frame}{Análisis causal}
\begin{alertblock}{Causal: ¿Qué causó este resultado?}
\begin{itemize}
    \item Identifica relaciones de \textbf{causa-efecto}
    \item Diferencia \textbf{correlación} de \textbf{causalidad}
\end{itemize}
\end{alertblock}

\begin{exampleblock}{Ejemplo}
``El aumento del precio \textbf{causó} una disminución del 12\% en la demanda''\\
(controlando por estacionalidad y competencia)
\end{exampleblock}

\vspace{0.5cm}
\begin{center}
\Large
\alert{Correlación $\neq$ Causalidad}
\end{center}
\end{frame}

\section{Roles en ciencia de datos}

\begin{frame}{Roles en ciencia de datos}
\begin{itemize}
    \item \textbf{Analista de datos}
    \begin{itemize}
        \item Análisis descriptivo y exploratorio
        \item Reportes y dashboards
        \item SQL, Excel, herramientas de BI
    \end{itemize}
\end{itemize}
\end{frame}

\begin{frame}{Roles en ciencia de datos}
\begin{itemize}
    \item \textbf{Analista de datos}
    \begin{itemize}
        \item Análisis descriptivo y exploratorio
        \item Reportes y dashboards
        \item SQL, Excel, herramientas de BI
    \end{itemize}
    
    \item \textbf{Científico de datos}
    \begin{itemize}
        \item Modelos predictivos, machine learning
        \item Experimentación, A/B testing
        \item Python, R, estadística avanzada
    \end{itemize}
\end{itemize}
\end{frame}

\begin{frame}{Roles en ciencia de datos}
\begin{itemize}
    \item \textbf{Analista de datos}
    \begin{itemize}
        \item Análisis descriptivo y exploratorio
        \item Reportes y dashboards
        \item SQL, Excel, herramientas de BI
    \end{itemize}
    
    \item \textbf{Científico de datos}
    \begin{itemize}
        \item Modelos predictivos, machine learning
        \item Experimentación, A/B testing
        \item Python, R, estadística avanzada
    \end{itemize}
    
    \item \textbf{Ingeniero de datos}
    \begin{itemize}
        \item Infraestructura, pipelines, ETL
        \item Bases de datos, data warehouses
        \item SQL, Spark, Airflow
    \end{itemize}
\end{itemize}
\end{frame}

\begin{frame}{Roles en ciencia de datos}
\begin{itemize}
    \item \textbf{Analista de datos}
    \begin{itemize}
        \item Análisis descriptivo y exploratorio
        \item Reportes y dashboards
        \item SQL, Excel, herramientas de BI
    \end{itemize}
    
    \item \textbf{Científico de datos}
    \begin{itemize}
        \item Modelos predictivos, machine learning
        \item Experimentación, A/B testing
        \item Python, R, estadística avanzada
    \end{itemize}
    
    \item \textbf{Ingeniero de datos}
    \begin{itemize}
        \item Infraestructura, pipelines, ETL
        \item Bases de datos, data warehouses
        \item SQL, Spark, Airflow
    \end{itemize}
    
    \item \textbf{Ingeniero de ML}
    \begin{itemize}
        \item Modelos en producción
        \item Escalabilidad, monitoring
        \item MLOps, deployment
    \end{itemize}
\end{itemize}
\end{frame}

\section{Configuración del entorno}

\begin{frame}{Herramientas del curso}
\begin{block}{Python + Anaconda}
\begin{itemize}
    \item Descarga: \texttt{anaconda.com}
    \item Incluye: pandas, numpy, scikit-learn, matplotlib, seaborn
    \item Gestores: \texttt{conda} y \texttt{pip}
\end{itemize}
\end{block}

\begin{block}{R + RStudio}
\begin{itemize}
    \item R: \texttt{cran.r-project.org}
    \item RStudio: \texttt{posit.co}
    \item Librerías: tidyverse, ggplot2, caret
\end{itemize}
\end{block}

\begin{block}{Visual Studio Code}
\begin{itemize}
    \item Descarga: \texttt{code.visualstudio.com}
    \item Extensiones: Python, Jupyter, R
    \item Interfaz principal para notebooks
\end{itemize}
\end{block}
\end{frame}

\begin{frame}{Git y GitHub}
\begin{block}{Git}
\begin{itemize}
    \item Control de versiones local
    \item Descarga: \texttt{git-scm.com}
    \item Rastrea cambios, colaboración
\end{itemize}
\end{block}

\begin{block}{GitHub}
\begin{itemize}
    \item Plataforma web para repositorios
    \item Crear cuenta: \texttt{github.com}
    \item Alojar código, colaborar, compartir
\end{itemize}
\end{block}

\vspace{0.5em}

\begin{alertblock}{Tarea para hoy}
\begin{enumerate}
    \item Instalar Anaconda (Python)
    \item Instalar R y RStudio
    \item Instalar VS Code + extensiones
\end{enumerate}
\end{alertblock}
\end{frame}

\section{Resumen}

\begin{frame}{Resumen de la clase}
\begin{itemize}
    \item El análisis de datos sigue un \textbf{ciclo iterativo}
    \item Diferentes tipos: descriptivo, exploratorio, inferencial, predictivo, causal
    \item Diferentes roles: analista, científico, ingeniero de datos/ML
    \item Herramientas: \textbf{Python, R, VS Code, Git, GitHub}
\end{itemize}

\vspace{1em}

\begin{block}{Próxima clase}
\begin{itemize}
    \item Introducción a Jupyter notebooks
    \item Sintaxis básica de Python y R
    \item Buenas prácticas y reproducibilidad
    \item Caso de estudio: dataset Iris
\end{itemize}
\end{block}
\end{frame}


\end{document}
